\documentclass[11pt,a4paper]{article}

% --- Accessibility & UX Packages ---
\usepackage[T1]{fontenc}
\usepackage[utf8]{inputenc}
\usepackage{lmodern} % Better font for screen readers
\usepackage{xcolor}  % Color contrast
\usepackage{amsmath} % Math formulas
\usepackage{amssymb}
\usepackage{geometry}
\geometry{margin=1in} % Good whitespace for readability

% Accessible hyperlinks
\definecolor{linkblue}{RGB}{0,86,179} % High contrast link color (WCAG AAA for text)
\usepackage[
    colorlinks=true,
    linkcolor=linkblue,
    citecolor=linkblue,
    urlcolor=linkblue,
    pdftitle={Geotemporal Hydrodynamics: The Unified Fluid Theory},
    pdfauthor={T. Abram},
    pdfsubject={Physics, Cosmology},
    pdfkeywords={Fluid Theory, Gravity, Cosmology},
    bookmarks=true,
    pdfdisplaydoctitle=true,
]{hyperref}

% --- Document Metadata ---
\title{Geotemporal Hydrodynamics\\ \Large The Unified Fluid Theory}
\author{T. Abram}
\date{March 2026}

\begin{document}

\maketitle

\begin{abstract}
Geotemporal Hydrodynamics (GTH) proposes that spacetime is a macroscopic quantum superfluid
rather than an empty geometric manifold. The framework is motivated by the vacuum-energy inconsistency between Quantum Field Theory and standard cosmology, and models the vacuum as a 5D
viscoelastic substrate with order parameter $\Psi = \rho e^{i\theta}$. In this picture, particles are not point objects
but stable topological defects (GeoKnots), while gravity emerges as the long-wavelength acoustic
response of the substrate.

Starting from a covariant hydrodynamic action augmented by Chern--Simons and Wess--Zumino--
Witten sectors, the manuscript derives coupled continuity--Euler dynamics, an effective acoustic metric,
and a constitutive map to a Poisson-limit gravitational equation. A closed parameterization is presented
for effective gravitational coupling and defect-energy mass scaling, including explicit dependence on
substrate quantities ($m_\psi, \rho_0, c_s, \tau, \eta$). At galactic scales, the theory predicts wake-mediated transport
and inverse-cascade contributions to rotation support; at compact-object scales, it predicts frequencydependent viscoelastic suppression consistent with ringdown-fidelity constraints.

The theory is framed as falsifiable through three near-term channels: rotation-curve residual tests
(SPARC), ringdown dispersion/damping constraints (LIGO/Virgo/KAGRA), and analog superfluidvortex experiments that probe topological stability and transport signatures.
\end{abstract}

\section*{Axiomatic Executive Summary}
To establish falsifiability and rigorous scale-invariance goals, this treatise develops a mathematical
program for the following phenomena as testable hydrodynamic constructions requiring full derivation
and validation:

\begin{itemize}
    \item \textbf{I. The Grand Unification of Newton’s Constant (G)}\\
    Gravity is not a fundamental force, but an induced acoustic pressure coupling. A proposed
    effective expression for Newton’s Gravitational Constant is written in terms of the speed of sound
    ($c_s$), the reduced Planck constant ($\hbar$), and the squared mass of the substrate scalar mode ($m_\psi$):
    \begin{equation}
        G_{eff} = \frac{3\pi\hbar c_s}{4m_\psi^2}
    \end{equation}
    This recovers the correct $m^3 kg^{-1} s^{-2}$ dimensionality without auxiliary scaling parameters.

    \item \textbf{II. The Acoustic Derivation of Fundamental Mass ($m_e$)}\\
    Rest mass is not an intrinsic property; it is the localized Gross-Pitaevskii kinetic bending energy
    of the substrate. By modeling the Electron as a fundamental Trefoil Knot ($3_1$) bound by a
    Compton-resonant phase closure, the 0.511 MeV scale is modeled from the background fluid
    line-tension ($T_{vacuum}$):
    \begin{equation}
        m_e = \sqrt{\frac{\rho_0 h^3}{4\pi c^3 m_\psi^2}} \ln(\Lambda)
    \end{equation}

    \item \textbf{III. The Topological Derivation of the Pauli Exclusion Principle}\\
    Vortices in standard superfluids are bosonic. However, by augmenting the 5D Master Action with
    a Wess-Zumino-Witten (WZW) topological framing term, the Finkelstein-Rubinstein theorem
    can be applied.
\end{itemize}

\end{document}
